\documentclass[11pt,a4paper]{article}

\usepackage[french]{babel}
\usepackage[T1]{fontenc}
\usepackage[utf8]{inputenc}
\usepackage[a4paper,margin=2.2cm]{geometry}
\usepackage{xcolor}
\usepackage{tgadventor} % titres — pendant de Raleway
\renewcommand{\familydefault}{\sfdefault}
\usepackage{tgheros}    % corps de texte — pendant de Hanken Grotesk
\usepackage{tgcursor}   % monospace
\usepackage{tcolorbox}
\tcbuselibrary{raster,skins,breakable}
\usepackage{enumitem}
\usepackage{titlesec}
\usepackage{parskip}
\usepackage{hyperref}

% ---------- Palette (identique à l'appli) ----------
\definecolor{accent}{HTML}{DF5F4D}
\definecolor{accentdark}{HTML}{C94E3D}
\definecolor{accentsoft}{HTML}{FBEAE7}
\definecolor{good}{HTML}{3F7D4F}
\definecolor{goodsoft}{HTML}{E9F2EB}
\definecolor{warn}{HTML}{B3392A}
\definecolor{warnsoft}{HTML}{F7E9E7}
\definecolor{ink}{HTML}{221D17}
\definecolor{muted}{HTML}{8A7F6F}
\definecolor{border}{HTML}{ECDFC4}
\definecolor{cardbg}{HTML}{FBF0DB}

\color{ink}
\hypersetup{colorlinks=true, urlcolor=accentdark, linkcolor=accentdark}

% ---------- Titres de section ----------
\titleformat{\section}
  {\sffamily\fontfamily{qag}\selectfont\bfseries\Large\color{accent}}
  {}{0pt}{}
  [{\color{border}\titlerule[1.1pt]}]
\titlespacing*{\section}{0pt}{22pt}{10pt}

\titleformat{\subsection}
  {\sffamily\fontfamily{qag}\selectfont\bfseries\large\color{ink}}
  {}{0pt}{}
\titlespacing*{\subsection}{0pt}{14pt}{4pt}

% ---------- Petits utilitaires ----------
\newcommand{\eyebrow}[1]{{\sffamily\fontfamily{qag}\selectfont\bfseries\small\color{accentdark}\MakeUppercase{#1}}\par\vspace{2pt}}

\newtcbox{\flaggood}{on line, boxrule=0pt, colback=goodsoft, colframe=goodsoft,
  colupper=good, arc=8pt, boxsep=1pt, left=8pt, right=8pt, top=4pt, bottom=4pt,
  fontupper=\sffamily\fontfamily{qag}\selectfont\bfseries\small}
\newtcbox{\flagwarn}{on line, boxrule=0pt, colback=warnsoft, colframe=warnsoft,
  colupper=warn, arc=8pt, boxsep=1pt, left=8pt, right=8pt, top=4pt, bottom=4pt,
  fontupper=\sffamily\fontfamily{qag}\selectfont\bfseries\small}

\newtcolorbox{notebox}{
  colback=accentsoft, colframe=accent, boxrule=0.8pt, arc=6pt,
  left=12pt, right=12pt, top=8pt, bottom=8pt
}

\newtcolorbox{technicalbox}{
  colback=white, colframe=muted, boxrule=1.1pt, arc=6pt,
  frame style={dash pattern=on 5pt off 4pt},
  left=14pt, right=14pt, top=10pt, bottom=10pt
}

\newtcolorbox{sqlbox}{
  colback=ink, colframe=ink, arc=4pt,
  left=12pt, right=12pt, top=6pt, bottom=6pt,
  fontupper=\ttfamily\small\color{cardbg}
}

\newtcolorbox{rolecard}[1][]{
  colback=white, colframe=border, boxrule=0.8pt, arc=8pt,
  left=12pt, right=12pt, top=10pt, bottom=10pt,
  valign=top, #1
}

\setlist[itemize]{label={\color{accent}\textbullet}, leftmargin=1.4em, itemsep=2pt, topsep=2pt}
\setlist[enumerate]{leftmargin=1.6em, itemsep=4pt, topsep=2pt}

\pagestyle{empty}

\begin{document}

% ---------- En-tête ----------
\begin{center}
\begin{minipage}{\textwidth}
{\sffamily\fontfamily{qag}\selectfont\bfseries\huge\color{accent}Fonctionnement de l'appli Planning}\\[3pt]
{\sffamily\itshape\small\color{muted}Gestion des plannings hebdomadaires — CO-P1}
\end{minipage}
\end{center}
\vspace{2pt}
{\color{border}\hrule height 1.4pt}
\vspace{8pt}

\noindent L'appli remplace le fichier Excel des plannings hebdo. Une \textbf{semaine} = une
grille du lundi au dimanche. Chaque jour contient un ou plusieurs \textbf{créneaux} (trajets /
tournées). Les \textbf{référents} et les \textbf{bénévoles} sont affectés à chaque créneau ; les
bénévoles s'inscrivent eux-mêmes via un \textbf{lien public}, sans compte.

\vspace{4pt}
\noindent Sur téléphone, le menu principal passe dans un \textbf{tiroir} : le bouton à trois
barres en haut à droite ouvre un panneau depuis la droite (fond flouté derrière, on referme en
touchant à côté).

% ---------- Rôles ----------
\section{2 rôles}

\begin{tcbraster}[raster columns=2, raster equal height=rows, raster column skip=8pt]
\begin{rolecard}
{\ttfamily\footnotesize\color{muted}Gestion des plannings}\\[6pt]
{\sffamily\fontfamily{qag}\selectfont\bfseries Gestionnaire}\\[6pt]
{\small Crée, édite et archive les semaines. Ajoute / modifie / supprime les créneaux. Gère les
référents et les bénévoles. Partage le lien public, verrouille les inscriptions.}\\[6pt]
{\small Ne gère pas les comptes. Ne supprime pas définitivement une semaine (l'archivage suffit).}
\end{rolecard}
\begin{rolecard}[colframe=accent, boxrule=1.4pt]
{\ttfamily\footnotesize\color{muted}Administration}\\[6pt]
{\sffamily\fontfamily{qag}\selectfont\bfseries Admin}\\[6pt]
{\small Tout ce que fait un gestionnaire, plus : dans l'onglet \textbf{Utilisateurs}, création de
comptes, modification NOM / prénom / email / téléphone / rôle, définition d'un mot de passe,
suppression d'un compte.}\\[6pt]
{\small Peut aussi supprimer \textbf{définitivement} une semaine.}
\end{rolecard}
\end{tcbraster}

\begin{notebox}
Les \textbf{seuls comptes} sont ceux des personnes qui gèrent le planning (une petite équipe).
Les bénévoles n'ont pas de compte : ils passent par le lien public.

\vspace{4pt}
Après avoir créé le tout premier compte (page \texttt{/inscription}) et l'avoir passé en
\texttt{admin}, on \textbf{désactive l'inscription} dans Supabase : plus personne ne peut créer
de compte de l'extérieur. Ensuite, les comptes se créent uniquement depuis l'onglet
\textbf{Utilisateurs}. (Détails dans le \texttt{README}.)
\end{notebox}

\newpage
% ---------- Semaines ----------
\section{Semaines}
\eyebrow{Onglet Semaines}

La page d'accueil liste les semaines \textbf{actives} (en cours et à venir).

\begin{itemize}
  \item \textbf{+ Semaine courante} / \textbf{+ Semaine suivante} : crée la semaine
  correspondante. Les dates sont calées sur l'heure de \textbf{Paris} — le lundi de la semaine
  en cours est calculé automatiquement, pas besoin de le saisir.
  \item \textbf{Dupliquer} : crée une nouvelle semaine en recopiant tous les créneaux et les
  \textbf{référents} d'une semaine existante (les bénévoles ne sont pas recopiés). L'appli
  propose la \textbf{première semaine libre} ; on peut en choisir une autre. Pratique quand
  les trajets se répètent d'une semaine à l'autre.
  \item \textbf{Surnom} : le bouton « surnom » à côté du titre ajoute un surnom facultatif
  (« Semaine de la rentrée »…). Il s'affiche \textbf{en plus} des dates, partout, sans les
  remplacer.
  \item Le lien \textbf{« Voir les semaines archivées »} mène aux semaines passées.
\end{itemize}

\flaggood{Créer une semaine est réversible} \; (archivage, ou suppression définitive réservée
aux admins)

% ---------- Une semaine ----------
\section{Éditer une semaine}
\eyebrow{Trois affichages : Liste, Grille, Vertical}

Un bandeau \textbf{Liste / Grille / Vertical} en haut de la semaine bascule entre trois
affichages (le choix est mémorisé) :
\begin{itemize}
  \item \textbf{Grille} : vue classique, colonnes = jours, lignes = heures, les 7 jours côte à
  côte. Idéale sur grand écran ; \textbf{masquée automatiquement sur téléphone} (trop étroite
  pour rester lisible sur un petit écran).
  \item \textbf{Vertical} : même principe (axe des heures vertical, créneaux qui se
  chevauchent côte à côte) mais \textbf{un jour à la fois, pleine largeur} — pensé pour le
  portrait mobile.
  \item \textbf{Liste} : les créneaux de chaque jour, triés par heure et regroupés par
  \textbf{Matin / Après-midi / Soir}.
\end{itemize}

En vue \textbf{Liste} et \textbf{Vertical}, chaque jour est un \textbf{bloc repliable} : seul
\textbf{aujourd'hui} s'ouvre automatiquement au chargement, les autres jours restent repliés
(un tap sur l'en-tête les ouvre) — moins de défilement sur une semaine chargée. Le bouton
flottant \textbf{« Aujourd'hui »} ramène directement au jour courant.

Quel que soit l'affichage, un créneau est un \textbf{bloc coloré} qui met les \textbf{horaires
en gros, en premier}, puis l'intitulé, les référents, les bénévoles et les notes (détail complet
dans l'infobulle au survol) ; \textbf{cliquer un bloc ouvre son éditeur}. Un créneau dont
l'horaire n'est pas interprétable va dans la bande \textbf{« heure ? »}. Le bouton
\textbf{+ Créneau} est en \textbf{haut} de chaque colonne / section de jour. En vue Grille, le
bouton \textbf{« Colonnes - / + »} élargit ou resserre les colonnes. Une \textbf{semaine
archivée} passe en lecture seule (ni + Créneau, ni édition des blocs) : cliquez
\textbf{« Réactiver »} pour la modifier de nouveau.

\subsection*{Un créneau contient :}
\begin{itemize}
  \item le \textbf{trajet / destination} (ex. \emph{Rungis}, \emph{BAPIF Saclay}) — champ à
  autocomplétion sur les \textbf{lieux enregistrés}
  \item les \textbf{horaires} en texte libre (\texttt{8h-11h}, \texttt{11h30-14h}, \texttt{17h}…)
  \item le \textbf{lieu de départ} (ex. \emph{Censier}, \emph{BAPIF Arcueil}) — même
  autocomplétion
  \item le \textbf{véhicule} : \texttt{Ducato (Manuel)}, \texttt{Jumpy (auto)},
  \texttt{Europcar (auto)} ou \texttt{Sans véhicule}
  \item \textbf{au moins un référent} (obligatoire) — nom ; téléphone et lien vers un compte
  facultatifs
  \item le \textbf{nombre de places bénévoles} (2 par défaut)
  \item des \textbf{notes} (optionnel)
\end{itemize}

Les \textbf{bénévoles} se gèrent depuis l'éditeur du créneau (bouton \textbf{« Gérer les
bénévoles »}), ou par le lien public. L'éditeur porte aussi \textbf{« Supprimer le créneau »}.
Quand un nom (trajet ou départ) correspond exactement à un lieu enregistré, son
\textbf{adresse} s'affiche sous le champ.

\begin{notebox}
\textbf{Ducato} apparaît avec une pastille « boîte manuelle » : pensez au permis manuelle pour
la personne qui conduit. \textbf{Europcar} est signalé en vert.
\end{notebox}

% ---------- Lieux et couleurs ----------
\section{Lieux et couleurs}
\eyebrow{Onglet Lieux}

Une liste partagée de lieux, chacun avec une \textbf{adresse} et une \textbf{couleur}.

\begin{itemize}
  \item L'\textbf{adresse} est proposée à la saisie d'un créneau (autocomplétion + rappel sous
  le champ).
  \item La \textbf{couleur} devient le fond du bloc du créneau, à l'écran \textbf{comme au PDF} :
  celle du lieu de \textbf{destination} s'il est enregistré, sinon celle du \textbf{départ},
  sinon gris. La correspondance est \textbf{exacte} — enregistrez chaque lieu précis
  (« BAPIF Arcueil », « BAPIF Saclay »…). Quelques lieux d'exemple sont pré-remplis ;
  complétez leurs adresses.
\end{itemize}

% ---------- Conflits de véhicule ----------
\section{Conflit de véhicule}
\eyebrow{Alerte automatique}

Si un \textbf{même véhicule} (hors « Sans véhicule ») est affecté à \textbf{deux créneaux du
même jour dont les horaires se chevauchent}, un bandeau rouge le signale en haut de la semaine
et les créneaux concernés sont encadrés en rouge.

\vspace{4pt}
Si des horaires ne sont pas interprétables (texte trop libre), le conflit n'est pas détecté —
mieux vaut ne rien signaler que signaler à tort.

% ---------- Lien public ----------
\section{Lien public — inscription des bénévoles}
\eyebrow{Bouton « Lien public »}

Le bouton propose deux liens (+ QR code) :
\begin{itemize}
  \item \textbf{Lien permanent} \texttt{…/semaine} : montre \textbf{toujours la semaine en
  cours} et se met à jour tout seul chaque lundi. C'est celui à partager une fois pour toutes
  (groupe WhatsApp, affichage sur place).
  \item \textbf{Lien de cette semaine précise} \texttt{…/p/<code>} (code court de
  6 caractères) : pour diffuser une semaine à l'avance.
\end{itemize}

La page publique reprend la \textbf{même présentation} que la version gestion (mêmes blocs
colorés). Sans compte, un bénévole peut :
\begin{itemize}
  \item consulter toute la semaine dans les \textbf{mêmes trois vues} que côté gestion (Liste,
  Grille, Vertical) — Liste choisie par défaut sur téléphone, Grille masquée sur petit écran ;
  \item \textbf{toucher un créneau} pour ouvrir sa fiche détaillée (jour, horaires, départ +
  adresse, destination, véhicule, référents et bénévoles avec téléphones, places libres, notes) ;
  \item toucher \textbf{« + M'inscrire »} (sur le créneau ou dans sa fiche) s'il reste des
  places, saisir son \textbf{nom} et — recommandé — son \textbf{téléphone} ;
  \item se \textbf{retirer} d'un créneau où il s'est inscrit depuis ce navigateur.
\end{itemize}

Contrôles côté serveur : créneau complet, doublon de nom, semaine verrouillée -- l'inscription
est refusée avec un message clair.

\vspace{6pt}
\flagwarn{Verrouiller} \; Le bouton \textbf{« Verrouiller »} coupe toutes les inscriptions
publiques ; la semaine reste consultable. Une semaine archivée est aussi en lecture seule.

\newpage
% ---------- Export PDF ----------
\section{Export PDF}
\eyebrow{Bouton « Exporter en PDF »}

Le bouton \textbf{télécharge directement un PDF} (A4 paysage, une page) : en-tête avec le
\textbf{logo}, « Semaine du … » (+ le surnom), un filet aux couleurs de l'appli, puis le
semainier coloré. Le \textbf{même format} est utilisé partout — export individuel d'une semaine
\emph{et} export groupé des archives. \textbf{Seule différence} : le \textbf{QR code} du lien
public n'apparaît que pour les semaines \textbf{actives}, pas pour les archivées.

\vspace{4pt}
\textbf{Export groupé} : dans l'onglet Archives, cochez plusieurs semaines puis
\textbf{« Exporter la sélection »} — un PDF par semaine, réunis dans un \texttt{.zip} (une
seule semaine cochée $\rightarrow$ PDF direct).

% ---------- Archivage ----------
\section{Archivage}
\eyebrow{Onglet Archives}

\begin{itemize}
  \item \textbf{Automatique} : une semaine dont le dimanche est passé bascule en
  \textbf{archivée} (au plus tard au prochain chargement de la liste des semaines).
  \item \textbf{Manuel} : bouton \textbf{« Archiver »} sur une semaine ; \textbf{« Réactiver »}
  pour la remettre en actif et pouvoir la modifier.
\end{itemize}

Une semaine archivée est en \textbf{lecture seule} mais reste \textbf{consultable et
ré-exportable en PDF} depuis l'onglet Archives (cf. \emph{Export PDF} pour l'export groupé).

% ---------- Historique ----------
\section{Historique des modifications}
\eyebrow{En bas de chaque semaine}

Chaque ajout / modification / suppression de créneau ou d'affectation est journalisé avec
\textbf{qui} (nom du compte, ou « Lien public » pour une inscription bénévole), \textbf{quand},
et un court résumé. La section « Historique » est dépliable en bas de la page d'une semaine.

% ---------- Technique ----------
\section{Repères techniques}
\begin{technicalbox}
\textbf{Base} : Supabase (PostgreSQL + Auth + Edge Functions).
Tables : \texttt{profiles}, \texttt{plannings}, \texttt{activites}, \texttt{affectations},
\texttt{historique}, \texttt{lieux}.

\vspace{4pt}
\textbf{Sécurité} : RLS active partout. Le public (\texttt{anon}) n'a \emph{aucun} accès direct
aux tables : il passe uniquement par des fonctions \texttt{SECURITY DEFINER}
(\texttt{planning\_public}, \texttt{planning\_public\_courant}, \texttt{inscrire\_benevole},
\texttt{desinscrire\_benevole}) qui prennent le code du lien en paramètre.

\vspace{4pt}
\textbf{Comptes} : gérés par des Edge Functions à clé \texttt{service\_role}
(\texttt{admin-create-user}, \texttt{admin-set-password}, \texttt{admin-update-email},
\texttt{delete-user}), chacune vérifiant que l'appelant est \texttt{admin}.

\vspace{4pt}
\textbf{Archivage auto} : fonction \texttt{archiver\_semaines\_ecoulees()}, appelée à
l'ouverture de la liste, et — si \texttt{pg\_cron} est activé — une fois par nuit.
\end{technicalbox}

% ---------- Pied de page ----------
\vspace{18pt}
{\color{border}\hrule height 1.4pt}
\vspace{10pt}
\noindent\small Pour toute suggestion, demande ou problème, contacter\\
{\sffamily\fontfamily{qag}\selectfont\bfseries Samuel LEGER}.

\end{document}
